% questions.tex —— 2012 年·上海卷·文科 · 题目内容片段（无答案）
% 由 tools/generate_projects.py 依来源总账自 2012/2012_数学真题_上海文_试题.md 生成，勿手改（改源 MD 后重跑）

\gkpart{一、填空题}

\begin{question}
\gkti{1}计算：$\dfrac{3-\mathrm i}{1+\mathrm i}=$\gkblank{}。（$\mathrm i$ 为虚数单位）

\end{question}
\begin{question}
\gkti{2}若集合 $A=\{x\mid2x-1>0\}$，$B=\{x\mid|x|<1\}$，则 $A\cap B=$\gkblank{}。

\end{question}
\begin{question}
\gkti{3}函数 $f(x)=\begin{vmatrix}\sin x&2\\ -1&\cos x\end{vmatrix}$ 的最小正周期是\gkblank{}。

\end{question}
\begin{question}
\gkti{4}若 $\overrightarrow d=(2,1)$ 是直线 $l$ 的一个方向向量，则 $l$ 的倾斜角的大小为\gkblank{}。（结果用反三角函数值表示）

\end{question}
\begin{question}
\gkti{5}一个高为 $2$ 的圆柱，底面周长为 $2\pi$。该圆柱的表面积为\gkblank{}。

\end{question}
\begin{question}
\gkti{6}方程 $4^x-2^{x+1}-3=0$ 的解是\gkblank{}。

\end{question}
\begin{question}
\gkti{7}有一列正方体，棱长组成以 $1$ 为首项、$\dfrac12$ 为公比的等比数列，体积分别记为 $V_1$，$V_2$，$\cdots$，$V_n$，$\cdots$，则 $\lim\limits_{n\to\infty}(V_1+V_2+\cdots+V_n)=$\gkblank{}。

\end{question}
\begin{question}
\gkti{8}在 $\left(x-\dfrac1x\right)^6$ 的二项展开式中，常数项等于\gkblank{}。

\end{question}
\begin{question}
\gkti{9}已知 $y=f(x)$ 是奇函数。若 $g(x)=f(x)+2$ 且 $g(1)=1$，则 $g(-1)=$\gkblank{}。

\end{question}
\begin{question}
\gkti{10}满足约束条件 $|x|+2|y|\leqslant2$ 的目标函数 $z=y-x$ 的最小值是\gkblank{}。

\end{question}
\begin{question}
\gkti{11}三位同学参加跳高、跳远、铅球项目的比赛。若每人只选择一个项目，则有且仅有两人选择的项目相同的概率是\gkblank{}。（结果用最简分数表示）

\end{question}
\begin{question}
\gkti{12}在矩形 $ABCD$ 中，边 $AB$，$AD$ 的长分别为 $2$，$1$。若 $M$，$N$ 分别是 $BC$，$CD$ 上的点，且满足 $\dfrac{|BM|}{|BC|}=\dfrac{|CN|}{|CD|}$，则 $\overrightarrow{AM}\cdot\overrightarrow{AN}$ 的取值范围是\gkblank{}。

\end{question}
\begin{question}
\gkti{13}已知函数 $y=f(x)$ 的图象是折线段 $ABC$，其中 $A(0,0)$，$B\left(\dfrac12,1\right)$，$C(1,0)$。函数 $y=xf(x)$（$0\leqslant x\leqslant1$）的图象与 $x$ 轴围成的图形的面积为\gkblank{}。

\end{question}
\begin{question}
\gkti{14}已知 $f(x)=\dfrac1{1+x}$，各项均为正数的数列 $\{a_n\}$ 满足 $a_1=1$，$a_{n+2}=f(a_n)$。若 $a_{2010}=a_{2012}$，则 $a_{20}+a_{11}$ 的值是\gkblank{}。

\end{question}
\gkpart{二、选择题}

\begin{question}
\gkti{15}若 $1+\sqrt2\mathrm i$ 是关于 $x$ 的实系数方程 $x^2+bx+c=0$ 的一个复数根，则\mbox{（　　）}

（A）$b=2$，$c=3$

（B）$b=2$，$c=-1$

（C）$b=-2$，$c=-1$

（D）$b=-2$，$c=3$

\end{question}
\begin{question}
\gkti{16}对于常数 $m$，$n$，「$mn>0$」是「方程 $mx^2+ny^2=1$ 表示的曲线是椭圆」的\mbox{（　　）}

（A）充分不必要条件

（B）必要不充分条件

（C）充分必要条件

（D）既不充分也不必要条件

\end{question}
\begin{question}
\gkti{17}在 $\triangle ABC$ 中，若 $\sin^2A+\sin^2B<\sin^2C$，则 $\triangle ABC$ 的形状是\mbox{（　　）}

（A）钝角三角形

（B）直角三角形

（C）锐角三角形

（D）不能确定

\end{question}
\begin{question}
\gkti{18}若 $S_n=\sin\dfrac\pi7+\sin\dfrac{2\pi}7+\cdots+\sin\dfrac{n\pi}7$（$n\in\mathbf N^*$），则在 $S_1$，$S_2$，$\cdots$，$S_{100}$ 中，正数的个数是\mbox{（　　）}

（A）$16$

（B）$72$

（C）$86$

（D）$100$

\end{question}
\gkpart{三、解答题}

\Needspace{0.4\textheight}
\begin{question}
\gkti{19}如图，在三棱锥 $P-ABC$ 中，$PA\perp$ 底面 $ABC$，$D$ 是 $PC$ 的中点。已知 $\angle BAC=\dfrac\pi2$，$AB=2$，$AC=2\sqrt3$，$PA=2$。求：

（1） 三棱锥 $P-ABC$ 的体积；

（2） 异面直线 $BC$ 与 $AD$ 所成角的大小（结果用反三角函数值表示）。
\nopagebreak[4]
\begin{center}% 2012 年上海卷文科第 19 题：三棱锥 P-ABC，PA⊥底面 ABC，D 为 PC 中点。
% 来源 PDF：普通高考/2012/2012上海文.pdf，第 1 页（印面 769）第 2 栏。
% 重绘依据：原图为斜二测式透视，A 在左前、B 在左后、C 在右，P 在 A 的正上方；
%   实线为 PB、PC、AB、BC、AC，虚线为 PA（被遮挡）与 AD。
\begin{tikzpicture}[scale=0.95, semithick, every node/.style={inner sep=1.5pt}]
  \coordinate (A) at (0.2,0.3);
  \coordinate (B) at (-0.9,-1.1);
  \coordinate (C) at (3.6,-0.4);
  \coordinate (P) at (0.2,3.4);
  \coordinate (D) at ($(P)!0.5!(C)$);
  \draw (P) -- (B);
  \draw (P) -- (C);
  \draw (A) -- (B);
  \draw (B) -- (C);
  \draw (A) -- (C);
  \draw[dashed] (P) -- (A);
  \draw[dashed] (A) -- (D);
  \node[left] at (A) {$A$};
  \node[below left] at (B) {$B$};
  \node[right] at (C) {$C$};
  \node[above] at (P) {$P$};
  \node[right] at ([xshift=2pt]D) {$D$};
\end{tikzpicture}
\end{center}
\nopagebreak[4]
\end{question}
\begin{question}
\gkti{20}已知函数 $f(x)=\lg(x+1)$。

（1） 若 $0<f(1-2x)-f(x)<1$，求 $x$ 的取值范围；

（2） 若 $g(x)$ 是以 $2$ 为周期的偶函数，且当 $0\leqslant x\leqslant1$ 时，有 $g(x)=f(x)$，求函数 $y=g(x)$（$x\in[1,2]$）的反函数。

\Needspace{0.4\textheight}
\end{question}
\begin{question}
\gkti{21}海事救援船对一艘失事船进行定位：以失事船的当前位置为原点，以正北方向为 $y$ 轴正方向建立平面直角坐标系（以 $1$ 海里为单位长度），则救援船恰好在失事船正南方向 $12$ 海里 $A$ 处，如图。现假设：①失事船的移动路径可视为抛物线 $y=\dfrac{12}{49}x^2$；②定位后救援船即刻沿直线匀速前往救援；③救援船出发 $t$ 小时后，失事船所在位置的横坐标为 $7t$。

（1） 当 $t=0.5$ 时，写出失事船所在位置 $P$ 的纵坐标。若此时两船恰好会合，求救援船速度的大小和方向；

（2） 问救援船的时速至少是多少海里才能追上失事船？
\nopagebreak[4]
\begin{center}% 2012 年上海卷文科第 21 题：失事船路径（抛物线 y=12x^2/49）、救援船出发点 A(0,-12)。
% 来源 PDF：普通高考/2012/2012上海文.pdf，第 2 页（印面 770）第 1 栏。
% 重绘依据：原图为平面直角坐标系：y 轴向上经过 O、A，A 在负 y 轴上；抛物线开口向上过 O，
%   右支上标点 P，A 与 P 之间连一直线，x 轴正向朝右。
\begin{tikzpicture}[scale=0.33, semithick, >=Stealth]
  % 坐标轴
  \draw[->] (-4.6,0) -- (5.0,0) node[below] {$x$};
  \draw[->] (0,-14.6) -- (0,6.0) node[left] {$y$};
  \node[below left] at (0,0) {$O$};
  % 抛物线 y = 12x^2/49（画到 |x|<=4.6）
  \draw[domain=-4.6:4.6, samples=120] plot (\x, {12*\x*\x/49});
  % 关键点
  \coordinate (A) at (0,-12);
  \coordinate (P) at (3.5,3);
  \fill (A) circle (2pt) node[left] {$A$};
  \fill (P) circle (2pt) node[right] {$P$};
  \draw (A) -- (P);
\end{tikzpicture}
\end{center}
\nopagebreak[4]
\end{question}
\begin{question}
\gkti{22}在平面直角坐标系 $xOy$ 中，已知双曲线 $C:2x^2-y^2=1$。

（1） 设 $F$ 是 $C$ 的左焦点，$M$ 是 $C$ 右支上一点。若 $|MF|=2\sqrt2$，求点 $M$ 的坐标；

（2） 过 $C$ 的左顶点作 $C$ 的两条渐近线的平行线，求这两组平行线围成的平行四边形的面积；

（3） 设斜率为 $k$（$|k|<\sqrt2$）的直线 $l$ 交 $C$ 于 $P$，$Q$ 两点。若 $l$ 与圆 $x^2+y^2=1$ 相切，求证：$OP\perp OQ$。

\end{question}
\begin{question}
\gkti{23}对于项数为 $m$ 的有穷数列 $\{a_n\}$，记 $b_k=\max\{a_1,a_2,\cdots,a_k\}$（$k=1,2,\cdots,m$），即 $b_k$ 为 $a_1$，$a_2$，$\cdots$，$a_k$ 中的最大值，并称数列 $\{b_n\}$ 是 $\{a_n\}$ 的控制数列。如 $1$，$3$，$2$，$5$，$5$ 的控制数列是 $1$，$3$，$3$，$5$，$5$。

（1） 若各项均为正整数的数列 $\{a_n\}$ 的控制数列为 $2$，$3$，$4$，$5$，$5$，写出所有的 $\{a_n\}$；

（2） 设 $\{b_n\}$ 是 $\{a_n\}$ 的控制数列，满足 $a_k+b_{m-k+1}=C$（$C$ 为常数，$k=1,2,\cdots,m$）。求证：$b_k=a_k$（$k=1,2,\cdots,m$）；

（3） 设 $m=100$，常数 $a\in\left(\dfrac12,1\right)$。若 $a_n=an^2-(-1)^{\frac{n(n+1)}2}n$，$\{b_n\}$ 是 $\{a_n\}$ 的控制数列，求 $(b_1-a_1)+(b_2-a_2)+\cdots+(b_{100}-a_{100})$。

\clearpage

\end{question}